\documentclass[12pt,a4paper]{article}
\usepackage{amsmath, amssymb, amsthm}
\usepackage{geometry}
\geometry{margin=1in}
\usepackage{hyperref}
\usepackage{enumitem}
\setlist{nosep}

\title{Stochastic Calculus: Concise Summary}
\author{Navid}
\date{\today}

\begin{document}

\setlength{\parindent}{0pt}

\maketitle
\tableofcontents

%-------------------
\section*{Preface}
This document provides a highly condensed summary of the main concepts in Stochastic Calculus, suitable for quick revision.

%-------------------
\section{Probability Space}



Definition: \textbf{Sample space}

$\Omega = \{\omega_i : i \in \mathcal{I}\}$ is the set of all possible outcomes of a random experiment.


Definition: \textbf{Event}

A subset of $\Omega$.


Definition: \textbf{Random variable}

$X : \Omega \to \mathbb{R}$ is a function mapping the sample space to real numbers.


Definition: \textbf{Probability measure}

$\mathbb{P}$ on $\Omega$ is defined by:
\begin{itemize}
	\item P1: $\mathbb{P}(\Omega) = 1$
	\item P2: $0 \leq \mathbb{P}(A) \leq 1$ for any event $A$
	\item P3: For disjoint $A_i$, $\mathbb{P}\left(\bigcup_{i \geq 1} A_i\right) = \sum_{i \geq 1} \mathbb{P}(A_i)$
\end{itemize}

Theorem: \textbf{Probability measure on finite $\Omega$}

If $\Omega $ is finite, the probability measure axioms are equivalent to:
\begin{itemize}
	\item P1*: $\mathbb{P}(\omega_i) \geq 0$ for all $i$
	\item P2*: For any event $A$, $\mathbb{P}(A) = \sum_{\omega \in A} \mathbb{P}(\omega)$
	\item P3*: $\sum_{i=1}^n \mathbb{P}(\omega_i) = 1$
\end{itemize}



Theorem: \textbf{Properties of probability measures}
\begin{itemize}
	\item P4: $\mathbb{P}(A^c) = 1 - \mathbb{P}(A)$
	\item P5: $\mathbb{P}(\varnothing) = 0$
	\item P6: $\mathbb{P}(A \cup B) = \mathbb{P}(A) + \mathbb{P}(B) - \mathbb{P}(A \cap B)$
	\item P7: If $A \subseteq B$, then $\mathbb{P}(A) \leq \mathbb{P}(B)$
\end{itemize}

Definition: \textbf{Distribution(law) of a random variable}

The distribution (law) of a random variable $X$ is given by its cumulative distribution function $F_X : \mathbb{R} \to [0,1]$,

$F_X(x) = \mathbb{P}(\{\omega \in \Omega \mid X(\omega) \leq x\})$.

Definition: \textbf{Distribution(law) for finite $\Omega$}

If $\Omega$ is finite, the distribution of $X$ is also given by its probability mass function $f_X : \mathbb{R} \to [0,1]$,

$f_X(x) = \mathbb{P}(\{\omega \in \Omega \mid X(\omega) = x\})$.

\vspace{1em}

Definition: \textbf{Equality and equality in distribution}
\vspace{0.5em}
Random variables $X$ and $Y$ are equal if $X(\omega) = Y(\omega)$ for all $\omega \in \Omega$.
They are equal in distribution if they have the same distribution.
\vspace{0.5em}
\begin{itemize}
	\item If two random variables are equal, then they are equal in distribution.
	\item Two random variables may be equal in distribution but not equal.
	\item Two random variables may be equal in distribution under one probability measure but not under another.
\end{itemize}
\vspace{1em}


Definition: \textbf{Sigma-algebra}
\vspace{0.5em}
A $\sigma$-algebra $\mathcal{F}$ of $\Omega$ is a subset of events such that:
\begin{itemize}
	\item[T1:] $\Omega \in \mathcal{F}$
	\item[T2:] If $A \in \mathcal{F}$ then $A^c \in \mathcal{F}$
	\item[T3:] If $A_1, A_2, \ldots \in \mathcal{F}$ then $\bigcup_{i \geq 1} A_i \in \mathcal{F}$
\end{itemize}


\vspace{1em}

Definition: \textbf{Measurable space}
\vspace{0.5em}
The pair $(\Omega, \mathcal{F})$ of a sample space and a $\sigma$-algebra is called a measurable space.
\vspace{1em}

Definition: \textbf{Finite partition}
\vspace{0.5em}
A family $\mathcal{P} = \{A_1, \ldots, A_n\}$ of events in $\Omega$ is a finite partition if:
\begin{itemize}
	\item[PF1:] $A_i \neq \varnothing$ for all $i$
	\item[PF2:] $A_i \cap A_j = \varnothing$ for $i \neq j$
	\item[PF3:] $\bigcup_{i=1}^n A_i = \Omega$
\end{itemize}
\vspace{1em}


Definition: \textbf{Sigma-algebra generated by a partition} : $\sigma$-algebra $\mathcal{F}$ is generated from a finite partition $\mathcal{P}$ if it is the smallest $\sigma$-algebra containing all elements of $\mathcal{P}$. $\mathcal{F}$ is denoted $\sigma(\mathcal{P})$, and the elements of $\mathcal{P}$ are the atoms of $\mathcal{F}$.
\vspace{1em}

Definition: \textbf{Sigma-algebra generated by a random variable}
The smallest $\sigma$-algebra $\mathcal{F}$ that makes $X$ a random variable on $(\Omega, \mathcal{F})$ is called the $\sigma$-algebra generated by $X$, denoted $\sigma(X)$.
\vspace{0.5em}
If $\operatorname{Card}(\Omega) < \infty$, $X$ takes values $x_1, \ldots, x_m$ and $A_i = \{\omega \in \Omega \mid X(\omega) = x_i\}$, then $\mathcal{P} = \{A_1, \ldots, A_m\}$ is a finite partition and $\sigma(X) = \sigma(\mathcal{P})$.

\vspace{1em}
Definition: \textbf{Random variable}
A random variable $X$ constructed on the measurable space $(\Omega, \mathcal{F})$ is a real-valued function (maps the sample space to real numbers) $X : \Omega \to \mathbb{R}$ such that for all $x \in \mathbb{R}$ : $\{\omega \in \Omega : X(\omega) \leq x\} \in \mathcal{F}$.

if $card(\Omega) < \infty$, then the expression can also be written as $\{\omega \in \Omega : X(\omega) = x\} \in \mathcal{F}$. meaning that $X$ is a random variable on the measurable space $(\Omega, \mathcal{F})$ if and only if the events that characterize its distribution are elements of $\mathcal{F}$ (\textbf{$X$ to be $\mathcal{F}$-measurable}).
\vspace{1em}

Theorem: \textbf{measurability:} Let $(\Omega, \mathcal{F})$ be a measurable space with $\operatorname{Card}(\Omega) < \infty$ and $\mathcal{P} = \{A_1, \ldots, A_n\}$ the finite partition generating $\mathcal{F}$. The function $X : \Omega \to \mathbb{R}$ is a random variable on that space ($X$ is $\mathcal{F}$-measurable) if and only if $X$ is constant on the atoms of $\mathcal{F}$.
\vspace{1em}


Theorem: \textbf{Linear combinations of measurable random variables}
If $\operatorname{Card}(\Omega) < \infty$ and $X, Y : \Omega \to \mathbb{R}$ are $\mathcal{F}$-measurable, then for all $a, b \in \mathbb{R}$, $aX + bY$ is also $\mathcal{F}$-measurable. Any linear combination of random variables on the same measurable space is a random variable of that space.
\vspace{1em}


Definition: \textbf{Probability measure}
Let $(\Omega, \mathcal{F})$ be a measurable space. A function $\mathbb{P} : \mathcal{F} \to [0,1]$ is a probability measure if:
\begin{itemize}
	\item[P1:] $\mathbb{P}(\Omega) = 1$
	\item[P2:] $0 \leq \mathbb{P}(A) \leq 1$ for all $A \in \mathcal{F}$
	\item[P3:] For disjoint $A_i \in \mathcal{F}$, $\mathbb{P}(\bigcup_{i \geq 1} A_i) = \sum_{i \geq 1} \mathbb{P}(A_i)$
\end{itemize}
\vspace{1em}

Definition: \textbf{Probability space}
The triple $(\Omega, \mathcal{F}, \mathbb{P})$ of a sample space, a $\sigma$-algebra, and a probability measure is called a probability space.
\vspace{1em}




%-------------------
\section{Stochastic Processes}

Definition: \textbf{Stochastic process}
\vspace{0.5em}
Let $(\Omega, \mathcal{F})$ be a measurable space. A stochastic process $X = \{X_t : t \in \mathcal{T}\}$ is a family of random variables, all built on the same measurable space $(\Omega, \mathcal{F})$, where $\mathcal{T}$ is a set of indices.
\vspace{1em}

Definition: \textbf{Filtration}
\vspace{0.5em}
A family $\mathbb{F} = \{\mathcal{F}_t : t \in \mathcal{T}\}$ of $\sigma$-algebras on $\Omega$ is a filtration if:
\begin{itemize}
	\item[F1:] $\mathcal{F}_t \subseteq \mathcal{F}$ for all $t$
	\item[F2:] If $t_1 \leq t_2$, then $\mathcal{F}_{t_1} \subseteq \mathcal{F}_{t_2}$
\end{itemize}
\vspace{1em}

Definition: \textbf{Adapted process}
\vspace{0.5em}
A stochastic process $X = \{X_t : t \in \mathcal{T}\}$ is adapted to the filtration $\mathbb{F} = \{\mathcal{F}_t : t \in \mathcal{T}\}$ if $X_t$ is $\mathcal{F}_t$-measurable for all $t$.
\vspace{1em}

Definition: \textbf{Generated filtration}
\vspace{0.5em}
The filtration $\mathbb{F} = \{\mathcal{F}_t : t \in \mathcal{T}\}$ is generated by the process $X = \{X_t : t \in \mathcal{T}\}$ if $\mathcal{F}_t = \sigma(\{X_s : s \leq t\})$ for all $t$.
\vspace{1em}


\vspace{1em}

Definition: \textbf{Stopping time}
\vspace{0.5em}
Let $(\Omega, \mathcal{F})$ be a measurable space with $\operatorname{Card}(\Omega) < \infty$ and filtration $\mathbb{F} = \{\mathcal{F}_t : t \in \{0,1,\ldots\}\}$. A stopping time $\tau$ is an $(\Omega, \mathcal{F})$-random variable taking values in $\{0,1,\ldots\}$ such that $\{\omega \in \Omega : \tau(\omega) \leq t\} \in \mathcal{F}_t$ for all $t$. In discrete time, it is equivalent to check $\{\tau = t\}$ for each $t$.
Intuitively, the time $\tau$ when one makes a decision is a stopping time if the decision is made based on the information available at that time.
\vspace{1em}

Theorem: If $\tau_1$ and $\tau_2$ are stopping times, then $\tau_1 \wedge \tau_2 \equiv \min\{\tau_1, \tau_2\}$ and $\tau_1 \vee \tau_2 \equiv \max\{\tau_1, \tau_2\}$ are also stopping times.
\vspace{1em}

Definition: \textbf{First hitting time:} Let $(\Omega, \mathcal{F})$ be a measurable space with $\operatorname{Card}(\Omega) < \infty$ and filtration $\mathbb{F} = \{\mathcal{F}_t : t \in \{0,1,\ldots\}\}$. Let $X = \{X_t : t \in \{0,1,\ldots\}\}$ be a stochastic process adapted to $\mathbb{F}$, and $B \subseteq \mathbb{R}$. The first time $X$ enters $B$ is
$\tau_B(\omega) = \min\{t \in \{0,1,\ldots\} : X_t(\omega) \in B\}$.

If $X_t(\omega)$ never hits $B$, define $\tau_B(\omega) = \infty$.
\vspace{1em}

Theorem: The first hitting time $\tau_B$ is a stopping time.
\vspace{1em}





%-------------------
\section{Conditional Expectation}

Definition: \textbf{Independence of events:}
Let $(\Omega, \mathcal{F}, \mathbb{P})$ be a probability space. Two events $A$ and $B$ are independent if $\mathbb{P}(A \cap B) = \mathbb{P}(A)\mathbb{P}(B)$.
\vspace{1em}

Definition: \textbf{Independence of random variables:}
Random variables $X$ and $Y$ are independent if for all $A \in \mathcal{P}_X$ and $B \in \mathcal{P}_Y$, $\mathbb{P}(A \cap B) = \mathbb{P}(A)\mathbb{P}(B)$, where $\mathcal{P}_X$ and $\mathcal{P}_Y$ are finite partitions generating $\sigma(X)$ and $\sigma(Y)$. Intuitively, $X$ and $Y$ are independent if knowing one gives no information about the other.
\vspace{1em}

Definition: \textbf{Conditional probability:}
Let $(\Omega, \mathcal{F}, \mathbb{P})$ be a probability space. For any event $A \in \mathcal{F}$ with $\mathbb{P}(A) > 0$, the conditional probability given $A$ is
$\mathbb{P}(B|A) = \frac{\mathbb{P}(B \cap A)}{\mathbb{P}(A)}$ for all $B \in \mathcal{F}$.
\vspace{1em}


Theorem: \textbf{Independence via conditional probability}
\vspace{0.5em}
Let $X$ and $Y$ be random variables built on the probability space $(\Omega, \mathcal{F}, \mathbb{P})$, $\operatorname{Card}(\Omega) < \infty$. Let $\mathcal{P}_X = \{A_1, \ldots, A_m\}$ and $\mathcal{P}_Y = \{B_1, \ldots, B_n\}$ be finite partitions generating $\sigma(X)$ and $\sigma(Y)$. If

\[ \forall A \in \mathcal{P}_X \text{ such that } \mathbb{P}(A) > 0, \mathbb{P}(B|A) = \mathbb{P}(B), \forall B \in \mathcal{P}_Y \]
or
\[ \forall B \in \mathcal{P}_Y \text{ such that } \mathbb{P}(B) > 0, \mathbb{P}(A|B) = \mathbb{P}(A), \forall A \in \mathcal{P}_X \]
then $X$ and $Y$ are independent.
\vspace{1em}

Definition: \textbf{Expectation:}
Let $X$ be a random variable built on the probability space $(\Omega, \mathcal{F}, \mathbb{P})$ such that $\operatorname{Card}(\Omega) < \infty$. The expectation of $X$, denoted $\mathbb{E}^{\mathbb{P}}[X]$ is
\[
\mathbb{E}^{\mathbb{P}}[X] = \sum_{\omega \in \Omega} X(\omega) \mathbb{P}(\omega).
\]
\vspace{1em}

\textbf{Properties:}
\vspace{0.5em}
Let $X$ and $Y$ be two random variables built on the probability space $(\Omega, \mathcal{F}, \mathbb{P})$, $\operatorname{Card}(\Omega) < \infty$. If $a$ and $b$ represent real numbers, then
\begin{itemize}
	\item[E1] $\mathbb{E}^{\mathbb{P}}[aX + bY] = a\mathbb{E}^{\mathbb{P}}[X] + b\mathbb{E}^{\mathbb{P}}[Y]$;
	\item[E2] If $\forall \omega \in \Omega$, $X(\omega) \leq Y(\omega)$ then $\mathbb{E}^{\mathbb{P}}[X] \leq \mathbb{E}^{\mathbb{P}}[Y]$;
	\item[E3] In general, $\mathbb{E}^{\mathbb{P}}[XY] \neq \mathbb{E}^{\mathbb{P}}[X]\mathbb{E}^{\mathbb{P}}[Y]$;
	\item[E4] If $X$ and $Y$ are independent, then $\mathbb{E}^{\mathbb{P}}[XY] = \mathbb{E}^{\mathbb{P}}[X]\mathbb{E}^{\mathbb{P}}[Y]$.
\end{itemize}
\vspace{1em}

Definition: \textbf{kth moment}
\vspace{0.5em}
Let $X$ be a random variable built on the probability space $(\Omega, \mathcal{F}, \mathbb{P})$ such that $\operatorname{Card}(\Omega) < \infty$. The kth moment of $X$, denoted $\mathbb{E}^{\mathbb{P}}\left[X^k\right]$ is the expectation of the variable $X^k$:
\[
\mathbb{E}^{\mathbb{P}}\left[X^k\right] = \sum_{\omega \in \Omega} X^k(\omega) \mathbb{P}(\omega) = \sum_{i=1}^{n} x_i^k f_X(x_i)
\]
where $x_1 < \ldots < x_n$ are the values taken by $X$ and $f_X$ is its probability mass function.
\vspace{1em}

Definition: \textbf{Variance}
\vspace{0.5em}
Let $X$ be a random variable built on the probability space $(\Omega, \mathcal{F}, \mathbb{P})$ such that $\operatorname{Card}(\Omega) < \infty$. The variance of $X$, denoted $\operatorname{Var}^{\mathbb{P}}[X]$ is the expectation of the random variable $\left(X - \mathbb{E}^{\mathbb{P}}[X]\right)^2$:
\begin{align}
\operatorname{Var}^{\mathbb{P}}[X] &= \sum_{\omega \in \Omega} \left(X(\omega) - \mathbb{E}^{\mathbb{P}}[X]\right)^2 \mathbb{P}(\omega) \\
&= \sum_{i=1}^{n} \left(x_i - \mathbb{E}^{\mathbb{P}}[X]\right)^2 f_X(x_i)
\end{align}
where $x_1 < \ldots < x_n$ are the values taken by $X$ and $f_X$ is its probability mass function.
\vspace{1em}

\textbf{Properties:}
\begin{itemize}
	\item[V1] $\operatorname{Var}^{\mathbb{P}}[X] \ge 0$;
	\item[V2] $\operatorname{Var}^{\mathbb{P}}[X] = \mathbb{E}^{\mathbb{P}}[X^2] - \left(\mathbb{E}^{\mathbb{P}}[X]\right)^2$;
	\item[V3] $\forall a,b \in \mathbb{R},\; \operatorname{Var}^{\mathbb{P}}[aX + b] = a^2\operatorname{Var}^{\mathbb{P}}[X]$;
	\item[V4] If $X$ and $Y$ are independent, then $\operatorname{Var}^{\mathbb{P}}[X + Y] = \operatorname{Var}^{\mathbb{P}}[X] + \operatorname{Var}^{\mathbb{P}}[Y]$.
\end{itemize}
\vspace{1em}

%-------------------

Definition: \textbf{Covariance}
Let $X$ and $Y$ be two random variables built on the probability space $(\Omega, \mathcal{F}, \mathbb{P})$ such that $\operatorname{Card}(\Omega) < \infty$. The covariance of $X$ and $Y$, denoted $\operatorname{Cov}^{\mathbb{P}}[X,Y]$, is the expectation of the random variable $(X - \mathbb{E}^{\mathbb{P}}[X])(Y - \mathbb{E}^{\mathbb{P}}[Y])$:
\[
\operatorname{Cov}^{\mathbb{P}}[X,Y] = \sum_{\omega \in \Omega} \bigl(X(\omega) - \mathbb{E}^{\mathbb{P}}[X]\bigr)\bigl(Y(\omega) - \mathbb{E}^{\mathbb{P}}[Y]\bigr)\,\mathbb{P}(\omega).
\]
\vspace{1em}
\textbf{Properties:}
\begin{itemize}
	\item[C1] $\operatorname{Cov}^{\mathbb{P}}[X,Y] = \mathbb{E}^{\mathbb{P}}[XY] - \mathbb{E}^{\mathbb{P}}[X]\,\mathbb{E}^{\mathbb{P}}[Y]$;
	\item[C2] If $X$ and $Y$ are independent, then $\operatorname{Cov}^{\mathbb{P}}[X,Y] = 0$;
	\item[C3] $\forall a,b \in \mathbb{R},\; \operatorname{Cov}^{\mathbb{P}}[aX_1 + bX_2, Y] = a\operatorname{Cov}^{\mathbb{P}}[X_1,Y] + b\operatorname{Cov}^{\mathbb{P}}[X_2,Y]$;
	\item[C4] $\operatorname{Var}^{\mathbb{P}}[X + Y] = \operatorname{Var}^{\mathbb{P}}[X] + \operatorname{Var}^{\mathbb{P}}[Y] + 2\operatorname{Cov}^{\mathbb{P}}[X,Y]$.
\end{itemize}
\vspace{1em}

Definition: \textbf{Conditional expectation:}
\vspace{0.5em}
Let the random variable $X$ be built on the probability space $(\Omega, \mathcal{F}, \mathbb{P})$ and let $\mathcal{G} \subseteq \mathcal{F}$ be the $\sigma$-algebra generated by the finite partition $\mathcal{P} = \{A_1,\ldots,A_n\}$ satisfying $\forall i \in \{1,\ldots,n\},\; \mathbb{P}(A_i) > 0$. The conditional expectation of $X$ given $\mathcal{G}$, denoted $\mathbb{E}^{\mathbb{P}}[X\mid\mathcal{G}]$, is
\[
\mathbb{E}^{\mathbb{P}}[X\mid\mathcal{G}](\omega) = \sum_{i=1}^{n} \frac{\mathbb{I}_{A_i}(\omega)}{\mathbb{P}(A_i)} \sum_{\omega^* \in A_i} X(\omega^*)\,\mathbb{P}(\omega^*) = \sum_{i=1}^{n} \mathbb{I}_{A_i}(\omega) \sum_{\omega^* \in A_i} X(\omega^*)\,\mathbb{P}(\omega^*|A_i).
\]
where $\mathbb{I}_{A_i}:\Omega \to \{0,1\}$ is the indicator function
\[
\mathbb{I}_{A_i}(\omega) = \begin{cases} 1 & \text{if } \omega \in A_i,\\ 0 & \text{if } \omega \notin A_i. \end{cases}
\]
\textbf{Properties:}
\vspace{0.5em}
Let $\mathcal{G}$ be the $\sigma$-algebra generated by the finite partition $\mathcal{P}=\{A_1,\ldots,A_n\}$. Then for any random variables $X,Y$ built on $(\Omega,\mathcal{F},\mathbb{P})$:
\begin{itemize}
	\item[EC1] If $X$ is $\mathcal{G}$-measurable, then $\mathbb{E}^{\mathbb{P}}[X\mid\mathcal{G}] = X$;
	\item[EC2] If $\mathcal{G}_1 \subseteq \mathcal{G}_2$ are sigma-algebras, then $\mathbb{E}^{\mathbb{P}}\bigl[\mathbb{E}^{\mathbb{P}}[X\mid\mathcal{G}_1]\mid\mathcal{G}_2\bigr] = \mathbb{E}^{\mathbb{P}}[X\mid\mathcal{G}_1]$;
	\item[EC3] If $\mathcal{G}_1 \subseteq \mathcal{G}_2$ are sigma-algebras, then $\mathbb{E}^{\mathbb{P}}\bigl[\mathbb{E}^{\mathbb{P}}[X\mid\mathcal{G}_2]\mid\mathcal{G}_1\bigr] = \mathbb{E}^{\mathbb{P}}[X\mid\mathcal{G}_1]$;
	\item[EC4] $\mathbb{E}^{\mathbb{P}}[X\mid\{\varnothing,\Omega\}] = \mathbb{E}^{\mathbb{P}}[X]$;
	\item[EC5] $\mathbb{E}^{\mathbb{P}}\bigl[\mathbb{E}^{\mathbb{P}}[X\mid\mathcal{G}]\bigr] = \mathbb{E}^{\mathbb{P}}[X]$;
	\item[EC6] If $Y$ is $\mathcal{G}$-measurable, then $\mathbb{E}^{\mathbb{P}}[XY\mid\mathcal{G}] = Y\,\mathbb{E}^{\mathbb{P}}[X\mid\mathcal{G}]$;
	\item[EC7] If $X$ and $Y$ are independent, then $\mathbb{E}^{\mathbb{P}}[X\mid\sigma(Y)] = \mathbb{E}^{\mathbb{P}}[X]$;
	\item[EC8] $\forall a,b \in \mathbb{R},\; \mathbb{E}^{\mathbb{P}}[aX + bY\mid\mathcal{G}] = a\,\mathbb{E}^{\mathbb{P}}[X\mid\mathcal{G}] + b\,\mathbb{E}^{\mathbb{P}}[Y\mid\mathcal{G}]$.
\end{itemize}
\vspace{1em}



%-------------------
\section{Martingale}

Definition: \textbf{Discrete-time martingale}

On the filtered probability space $(\Omega, \mathcal{F}, \mathbb{F}, \mathbb{P})$, where $\mathbb{F}$ is the filtration $\{\mathcal{F}_t : t \in \{0, 1, 2, \ldots\}\}$, the stochastic process
$$M = \{M_t : t \in \{0, 1, 2, \ldots\}\}$$
is a discrete-time \textit{martingale} if:

\begin{itemize}
	\item (M1) $\forall t \in \{0, 1, 2, \ldots\}$, $\mathbb{E}^{\mathbb{P}}[|M_t|] < \infty$ \quad (being integrable)
	\item (M2) $\forall t \in \{0, 1, 2, \ldots\}$, $M_t$ is $\mathcal{F}_t$--measurable
	\item (M3) $\forall s, t \in \{0, 1, 2, \ldots\}$ such that $s < t$, $\mathbb{E}^{\mathbb{P}}[M_t \mid \mathcal{F}_s] = M_s$
\end{itemize}
\vspace{0.5em}

Lemma: \textbf{Equivalent third condition}

In the definition of a martingale, the condition (M3) is equivalent to
$$\text{M3*} \quad \forall t \in \{1, 2, \ldots\}, \quad \mathbb{E}^{\mathbb{P}}[M_t \mid \mathcal{F}_{t-1}] = M_{t-1}.$$
\vspace{0.5em}

Definition: \textbf{Stopped process}

The stochastic process $X$ and the stopping time $\tau$ are built on the same filtered measurable space $(\Omega, \mathcal{F}, \mathbb{F})$. The stochastic process $X^{\tau}$ defined by
$$X_t^{\tau}(\omega) = X_{t \wedge \tau(\omega)}(\omega)$$
is called a \textit{stopped process} with stopping time $\tau$.
\vspace{0.5em}

Theorem: \textbf{Stopped martingale}

If the martingale $M$ and the stopping time $\tau$ are built on the same filtered probability space $(\Omega, \mathcal{F}, \mathbb{F}, \mathbb{P})$, then the stopped process $M^{\tau}$ is also a martingale on that space.
\vspace{0.5em}

Theorem: \textbf{Optional Stopping Theorem}
\vspace{0.5em}

Let
$$X = \{X_t : t \in \{0, 1, 2, \ldots\}\}$$
be a process built on the filtered probability space $(\Omega, \mathcal{F}, \mathbb{F}, \mathbb{P})$, where $\mathbb{F}$ is the filtration $\{\mathcal{F}_t : t \in \{0, 1, 2, \ldots\}\}$. Let's assume that the stochastic process $X$ is $\mathbb{F}$--adapted and that it is integrable, i.e. $\mathbb{E}^{\mathbb{P}}[|X_t|] < \infty$. Then $X$ is a martingale if and only if
$$\mathbb{E}^{\mathbb{P}}[X_{\tau}] = \mathbb{E}^{\mathbb{P}}[X_0]$$
for any bounded stopping time $\tau$, i.e. for any given stopping time $\tau$, there exists a constant $b$ such that
$$\forall \omega \in \Omega, \quad 0 \leq \tau(\omega) < b.$$
\vspace{0.5em}

Definition: \textbf{Markovian process}
\vspace{0.5em}

A stochastic process $X = \{X_t : t \in \mathcal{T}\}$, where $\mathcal{T}$ is a set of indices, is said to be \textit{markovian} if, for any $t_1 < t_2 < \ldots < t_n \in \mathcal{T}$, the conditional distribution of $X_{t_n}$ given $X_{t_1}, \ldots, X_{t_{n-1}}$ is equal to the conditional distribution of $X_{t_n}$ given $X_{t_{n-1}}$, i.e. for any $x_1, \ldots, x_n \in \mathbb{R}$,
$$\mathbb{P}[X_{t_n} \leq x_n \mid X_{t_1} = x_1, \ldots, X_{t_{n-1}} = x_{n-1}] = \mathbb{P}[X_{t_n} \leq x_n \mid X_{t_{n-1}} = x_{n-1}].$$
\vspace{0.5em}

Definition: \textbf{Random walk}
\vspace{0.5em}

The stochastic process $X$, built on the probability space $(\Omega, \mathcal{F}, \mathbb{P})$, is a \textit{random walk} if it admits the representation
$$X_0 = 0 \quad \text{and} \quad \forall t \in \{1, 2, \ldots\}, \quad X_t = \sum_{n=1}^{t} \xi_n$$
where the sequence $\{\xi_i : i \in \{1, 2, \ldots\}\}$ is made up of independent and identically distributed random variables. Random walks are Markovian processes.
\vspace{0.5em}


%-------------------
\section{Discrete Time Market Model}

Riskless: $S_t^0 = (1+r)^t$, $S_0^0 = 1$. Risky: $S_t^1$ binary $\{s_1, s_2\}$, initial $s_0$. Portfolio $\phi = (\phi_0, \phi_1)$: $V_t(\phi) = \phi_0 S_t^0 + \phi_1 S_t^1$, cost $V_0(\phi) = \phi_0 + \phi_1 s_0$.
\vspace{1em}

Definition: \textbf{Arbitrage opportunity}
\vspace{0.5em}
A portfolio $\phi$ is an arbitrage opportunity if:
\begin{itemize}
	\item[A1:] $V_0(\phi) = 0$ (zero initial investment)
	\item[A2:] $\forall \omega \in \Omega$, $V_1(\omega) \geq 0$ (zero probability of losing money)
	\item[A3:] $\exists \omega \in \Omega$, $V_1(\omega) > 0$ (positive probability of gaining)
\end{itemize}
\vspace{1em}
So, arbitrage is possible if and only if $s_2 > s_1 \geq s_0(1+r)$ or $s_1 < s_2 \leq s_0(1+r)$.
We assume the model contains no arbitrage opportunities, i.e., $s_1 < s_0(1+r) < s_2$.
\vspace{1em}

Definition: \textbf{Contingent claim:}
A contingent claim is a contract (seller and buyer) whose value depends on the market state during the contract validity period (like insurance). Mathematically, a contingent claim $C$ is any non-negative random variable.
\vspace{1em}

\textbf{Replicating portfolio:}
 $\phi = (\phi_0, \phi_1)$ replicates contingent claim $C$ if $V_1(\phi) = C$. If replicating portfolio exists, $\phi_0 = \frac{s_0 c_1 - s_1 c_2}{(s_2-s_1)(1+r)}$ and $\phi_1 = \frac{c_2 - c_1}{s_2 - s_1}$.
\vspace{1em}

\textbf{Risk-neutral measure}
Probability measure $\mathbb{Q}$ defined by $q = \mathbb{Q}(\omega_1) = \frac{s_2 - s_0(1+r)}{s_2 - s_1}$ and $1-q = \mathbb{Q}(\omega_2)$. Under no-arbitrage, $0 < q < 1$.
\vspace{1em}

\textbf{Martingale property}
Under $\mathbb{Q}$: expected return on risky security = $r$ (same as riskless). Discounted price process $\frac{S_t^1}{(1+r)^t}$ is $\mathbb{Q}$-martingale.
\vspace{1em}

Definition: \textbf{Return on a security:}
The return on a security from $t-1$ to $t$ 
\[
R_t(\omega) = \frac{S_t(\omega) - S_{t-1}(\omega)}{S_{t-1}(\omega)}
\]
\vspace{1em}

Definition: \textbf{Price process}
\vspace{0.5em}
The $(K+1)$-dimensional stochastic process $\mathbf{S} = \{\mathbf{S}_t\}_{t \in \{0,1,\ldots,T\}}$ representing the evolution of individual unit prices of all securities is made up of random column vectors $\mathbf{S}_t = (S_t^0, S_t^1, \ldots, S_t^K)^\top$, where $S_t^k$ is the unit price of security $k$ at time $t$.
\vspace{1em}

Definition: \textbf{Riskless security}
The 0th security plays a peculiar role in that it is a riskless security (e.g., bank account or bond). It implies that $\forall \omega \in \Omega$ and $\forall t \in \{0, 1, \ldots, T-1\}$: $S_t^0(\omega) \leq S_{t+1}^0(\omega)$. Without loss of generality, $S_0^0(\omega) = 1$.
\vspace{1em}

Definition: \textbf{Discount factor}
The unidimensional stochastic process $\beta_t = \frac{1}{S_t^0}$ is used as discount factor.
\vspace{1em}

Definition: \textbf{Trading strategy:}
A trading strategy is a $\{\mathcal{F}_t\}_{t \in \{1, \ldots, T\}}$-predictable vector stochastic process $\phi = \{\phi_t\}_{t \in \{1, \ldots, T\}}$, where the random row vector $\phi_t = (\phi_t^0, \phi_t^1, \ldots, \phi_t^K)$ represents the investor's portfolio at time $t$. Here, $\phi_t^k$ is the number of shares of security $k$ held during $(t-1,t]$.
\vspace{1em}

Definition: \textbf{Self-financing investment strategy}
\vspace{0.5em}

A strategy is self-financing if no funds are added to or withdrawn from the value of the portfolio after time $t = 0$, i.e., $\forall t \in \{1, \ldots, T-1\}$: $\phi_t \mathbf{S}_t = \phi_{t+1} \mathbf{S}_t$. assuming zero transaction costs.
\vspace{1em}

Definition: \textbf{Market value}
\vspace{0.5em}

The process $V(\phi) = \{V_t(\phi) : t \in \{0, 1, \ldots, T\}\}$ represents the market value of the strategy at any time:
\[
V_t(\phi) = \begin{cases} \phi_1 \mathbf{S}_0 & \text{if } t = 0 \\ \phi_t \mathbf{S}_t & \text{if } t \in \{1, \ldots, T\} \end{cases}
\]
\vspace{1em}

Definition: \textbf{Admissible investment strategies}
\vspace{0.5em}

A trading strategy is called admissible if it is self-financing and its market value is never negative. The set $\Phi$ of admissible strategies is:
\[
\Phi = \left\{ \phi \,\bigg|\, \begin{array}{l} \text{each component } \phi^k \text{ of } \phi \text{ is predictable,} \\ \phi \text{ is self-financing,} \\ \text{and } \forall t \in \{0, \ldots, T\}: V_t(\phi) \geq 0 \end{array} \right\}
\]
\vspace{1em}

Definition: \textbf{Arbitrage opportunity}
\vspace{0.5em}

An admissible strategy $\phi$ is an arbitrage opportunity if $V_0(\phi) = 0$ and if $\mathbb{E}^{\mathbb{P}}[V_T(\phi)] > 0$.
\vspace{1em}

Definition: \textbf{Contingent claim}
\vspace{0.5em}

A contingent claim $X$ is a $(\Omega, \mathcal{F}_T)$-non-negative random variable. It can be thought of as a contract between two parties, such that $X(\omega)$ will be paid by one party to the other if state $\omega$ pertains. The set $\mathbf{X}$ of contingent claims is:
\[
\mathbf{X} = \left\{ X \,\bigg|\, \begin{array}{l} X \text{ is a } (\Omega, \mathcal{F}_T)\text{-random variable} \\ \text{such that } \forall \omega \in \Omega, X(\omega) \geq 0 \end{array} \right\}
\]
\vspace{1em}

Definition: \textbf{Attainable contingent claim}
\vspace{0.5em}

A contingent claim $X$ is said to be attainable if there exists an admissible trading strategy $\phi$ that can replicate the cash flow generated by $X$, i.e., $V_T(\phi) = X$.
\vspace{1em}

Definition: \textbf{Price system}
\vspace{0.5em}

A price system $\pi$ is a linear operator on $\mathbf{X}$ that returns non-negative values $\pi : \mathbf{X} \to [0, \infty)$ and that satisfies:
\[
\pi(X) = 0 \Leftrightarrow X = 0
\]
and
\[
\forall a, b \geq 0 \text{ and } \forall X_1, X_2 \in \mathbf{X}, \quad \pi(aX_1 + bX_2) = a\pi(X_1) + b\pi(X_2).
\]

\textit{Intuition:} We are implicitly assuming that we are small investors so if we buy more the price won't go up (orderbook effect) or buying more won't make the price less (like Costco effect :))
\vspace{1em}

Definition: \textbf{Consistency}
\vspace{0.5em}

A price system is said to be consistent with the market model if the price associated with the contingent claim $V_T(\phi)$ is its market value at time $t = 0$, $V_0(\phi)$, i.e. $\pi(V_T(\phi)) = V_0(\phi)$.

The set $\Pi$ of price systems consistent with the market model is:
\[
\Pi = \left\{ \pi : \mathbf{X} \to [0, \infty) \,\bigg|\, \begin{array}{l} \pi(X) = 0 \Leftrightarrow X = 0, \\[0.3em] \forall a, b \geq 0 \text{ and } \forall X_1, X_2 \in \mathbf{X}, \\[0.3em] \pi(aX_1 + bX_2) = a\pi(X_1) + b\pi(X_2), \\[0.3em] \forall \phi \in \Phi, \, \pi(V_T(\phi)) = V_0(\phi) \end{array} \right\}
\]
\vspace{1em}

Definition: \textbf{Equivalent martingale measures}
\vspace{0.5em}

Let $\mathbf{P}$ denote the set of martingale measures equivalent to $\mathbb{P}$:
\[
\mathbf{P} = \left\{ \mathbb{Q} \,\bigg|\, \begin{array}{l} \mathbb{Q} \text{ is a probability measure on } (\Omega, \mathcal{F}), \\[0.3em] \forall \omega \in \Omega, \mathbb{Q}(\omega) > 0, \text{ and} \\[0.3em] \forall k \in \{0, 1, \ldots, K\}, \, \beta S^k \text{ is a } \mathbb{Q}\text{-martingale} \end{array} \right\}
\]
\vspace{1em}

Definition: \textbf{Predictable process}
\vspace{0.5em}

A stochastic process $X = \{X_t : t = 0, 1, 2, \ldots\}$ built on a measurable space $(\Omega, \mathcal{F})$ is said to be predictable with respect to the filtration $\{\mathcal{F}_t : t = 0, 1, 2, \ldots\}$ if $X_0$ is a $(\Omega, \mathcal{F}_0)$-random variable and $\forall t \in \{1, 2, \ldots\}$, $X_t$ is a $(\Omega, \mathcal{F}_{t-1})$-random variable.
\vspace{1em}

Theorem: \textbf{Bijective correspondence (Proposition 2.6)}
\vspace{0.5em}

There is bijective correspondence between the set $\Pi$ of price systems that are consistent with the market model and the set $\mathbf{P}$ of martingale measures equivalent to $\mathbb{P}$. Such a correspondence is given by:
\begin{itemize}
	\item[(i)] $\pi(X) = \mathbb{E}^{\mathbb{Q}}[\beta_T X]$, $X \in \mathbf{X}$
	\item[(ii)] $\mathbb{Q}(A) = \pi(S_T^0 \mathbb{I}_A)$, $A \in \mathcal{F}$
\end{itemize}

Intuition: if there exists a probability measure under
which all discounted price processes are martingales, then we
can build, based on such a measure, a price system for
attainable contingent claims that is consistent with the market
model. and vice versa.
\vspace{1em}

Definition: \textbf{Equivalence of two probability measures}
\vspace{0.5em}

Let $\mathbb{P}$ and $\mathbb{Q}$ be two probability measures that exist on the measurable space $(\Omega, \mathcal{F})$. The measures $\mathbb{P}$ and $\mathbb{Q}$ are said to be equivalent if and only if the impossible events are the same under both measures, i.e.,
\[
\forall A \in \mathcal{F}, \quad \mathbb{P}(A) = 0 \Leftrightarrow \mathbb{Q}(A) = 0.
\]

In our case, since $\forall \omega \in \Omega$, $\mathbb{P}(\omega) > 0$, all measures $\mathbb{Q}$ equivalent to $\mathbb{P}$ shall satisfy the condition
\[
\forall \omega \in \Omega, \quad \mathbb{Q}(\omega) > 0.
\]
\vspace{1em}

Theorem: \textbf{Lemma on page 228}
\vspace{0.5em}

If there exists a self-financing strategy $\phi$ (not necessarily admissible) such that $V_0(\phi) = 0$, $V_T(\phi) \geq 0$ and $\mathbb{E}^{\mathbb{P}}[V_T(\phi)] > 0$, then there exists an arbitrage opportunity.

\textbf{Interpretation:}
\begin{itemize}
	\item This lemma provides us with a necessary and sufficient condition for the existence of an arbitrage opportunity.
	\item The advantage of this result is that, from now on, we don't have to verify any more whether the strategy that is a candidate for an arbitrage opportunity is admissible.
\end{itemize}
\vspace{1em}

Theorem: \textbf{Fundamental theorem of asset pricing (Theorem 2.7)}
\vspace{0.5em}

The market model contains no arbitrage opportunities if and only if there exists at least one martingale measure equivalent to $\mathbb{P}$.
\vspace{1em}

Corollary: \textbf{The law of one price}
\vspace{0.5em}

If the market model contains no arbitrage, then there is a single price associated with any attainable contingent claim $X$ and it satisfies $\pi = \mathbb{E}^{\mathbb{Q}}[\beta_T X]$ for each martingale measure $\mathbb{Q} \in \mathbf{P}$.
\vspace{1em}

Theorem: \textbf{Proposition 2.8}
\vspace{0.5em}

If $\phi$ is an admissible strategy, then the process $\beta V(\phi)$, representing its discounted market value, is a $\mathbb{Q}$-martingale for each measure $\mathbb{Q} \in \mathbf{P}$.
\vspace{1em}

Definition: \textbf{Complete market}
\vspace{0.5em}

A market is said to be complete if it contains no arbitrage opportunity and if all contingent claims are attainable.
\vspace{1em}

Theorem: \textbf{Characterization of complete markets}
\vspace{0.5em}

A market is complete if and only if there exists a unique martingale measure.
\vspace{1em}

%-------------------
\section{Brownian Motion}


Definition: \textbf{Standard Brownian motion}
\vspace{0.5em}

A \textit{standard Brownian motion} $\{W_t : t \geq 0\}$ is an adapted stochastic process, built on a filtered probability space $(\Omega, \mathcal{F}, \mathbb{F}, \mathbb{P})$ such that:

\begin{itemize}
	\item (MB1) $\forall \omega \in \Omega$, $W_0(\omega) = 0$
	\item (MB2) $\forall 0 \leq t_0 \leq t_1 \leq \ldots \leq t_k$, the random variables $W_{t_1} - W_{t_0}$, $W_{t_2} - W_{t_1}$, $\ldots$, $W_{t_k} - W_{t_{k-1}}$ are independent
	\item (MB3) $\forall s, t \geq 0$ such that $s < t$, the random variable $W_t - W_s$ is normally distributed with expectation 0 and variance $t - s$, i.e. $W_t - W_s \sim N(0, t - s)$
	\item (MB4) $\forall \omega \in \Omega$, the path $t \to W_t(\omega)$ is continuous
\end{itemize}
\vspace{1em}


\textbf{Technical condition} ( $(\Omega, \mathcal{F}, \mathbb{F}, \mathbb{P})$ is our filtered probability space.
)

Since we will work with almost-sure equalities, we require that the set of events which have a zero-probability to occur be included in the sigma-algebra $\mathcal{F}_0$, which is to say that the set
\[
\mathcal{N} = \{A \in \mathcal{F} : \mathbb{P}(A) = 0\} \subseteq \mathcal{F}_0.
\]

Thus, if $X$ is $\mathcal{F}_t$--measurable and $Y = X$ $\mathbb{P}$--almost-surely, then we know that $Y$ is $\mathcal{F}_t$--measurable.
\vspace{1em}

In general, the filtration we use is $\mathbb{F} = \{\mathcal{F}_t : t > 0\}$ where
\[
\mathcal{F}_t = \sigma(\{W_s : 0 \leq s \leq t\} \cup \mathcal{N})
\]
is the smallest sigma-algebra for which the random variables $W_s : 0 \leq s \leq t$ are measurable and that contains the zero-measure sets.
\vspace{1em}

\textbf{Reminder about the Gaussian distribution}

If $X \sim N(\mu, \sigma^2)$, then $f_X(x) = \frac{1}{\sigma\sqrt{2\pi}} \exp\left\{-\frac{(x-\mu)^2}{2\sigma^2}\right\}$, $F_X(x) = \int_{-\infty}^{x} f_X(y) dy$, and $\forall a, b \in \mathbb{R}, a < b$: $\mathbb{P}[a < X \leq b] = \int_a^b f_X(x) dx$.
\vspace{1em}

Theorem: \textbf{Independence and covariance for multivariate normal}
\vspace{0.5em}

If $X$ and $Y$ are two random variables following a multivariate normal distribution, then $X$ and $Y$ are independent if and only if their covariance is zero.
\vspace{1em}

\textit{Note:} In general, independence guarantees zero covariance, but the reverse is not necessarily true.
\vspace{1em}

Theorem: \textbf{Brownian motion is a martingale and a Markovian process}
\vspace{1em}

Definition: \textbf{Continuous-time martingale}
\vspace{0.5em}

On the filtered probability space $(\Omega, \mathcal{F}, \mathbb{F}, \mathbb{P})$, where $\mathbb{F}$ is the filtration $\{\mathcal{F}_t : t \geq 0\}$, the stochastic process $M = \{M_t : t \geq 0\}$ is a martingale in continuous time if:

\begin{itemize}
	\item (M1) $\forall t \geq 0$, $\mathbb{E}^{\mathbb{P}}[|M_t|] < \infty$
	\item (M2) $\forall t \geq 0$, $M_t$ is $\mathcal{F}_t$--measurable
	\item (M3) $\forall s, t \geq 0$ such that $s < t$, $\mathbb{E}^{\mathbb{P}}[M_t \mid \mathcal{F}_s] = M_s$
\end{itemize}
\vspace{1em}

Theorem: \textbf{Brownian paths are nowhere differentiable}
\vspace{0.5em}

$\forall \omega \in \Omega$, the path $t \to W_t(\omega)$ is nowhere differentiable.
\vspace{1em}

Definition: \textbf{Multidimensional Brownian motion}
\vspace{0.5em}

Standard Brownian motion $W$ of dimension $n$ is a family of random vectors
\[
\left\{ W_t = \left( W_t^{(1)}, \ldots, W_t^{(n)} \right)^{\top} : t \geq 0 \right\}
\]
where $W^{(1)}, \ldots, W^{(n)}$ represent independent Brownian motions built on the filtered probability space $(\Omega, \mathcal{F}, \mathbb{F}, \mathbb{P})$.
\vspace{1em}

Theorem: \textbf{Linear transformation of independent Brownian motions}
\vspace{0.5em}

Let $W = \left(W^{(1)}, \ldots, W^{(n)}\right)^{\top}$ be a vector of independent Brownian motions and $\Gamma = (\gamma_{ij})_{i,j \in \{1,2,\ldots,n\}}$ be a matrix of constants. For all $t$, define
\[
B_t = \Gamma W_t.
\]
Then $B_t$ is a random vector of dimension $n$, the $i$th component of which is $B_t^{(i)} = \sum_{k=1}^{n} \gamma_{ik} W_t^{(k)}$. Moreover
\[
\operatorname{Cov}\left[B_t^{(i)}, B_t^{(j)}\right] = t \sum_{k=1}^{n} \gamma_{ik}\gamma_{jk}
\]
and
\[
\operatorname{Cor}\left[B_t^{(i)}, B_t^{(j)}\right] = \frac{\sum_{k=1}^{n} \gamma_{ik}\gamma_{jk}}{\sqrt{\sum_{k=1}^{n} \gamma_{ik}^2} \sqrt{\sum_{k=1}^{n} \gamma_{jk}^2}}.
\]
\vspace{1em}

Theorem: \textbf{Representation of correlated Brownian motions}
\vspace{0.5em}

Let $B^{(1)}, \ldots, B^{(n)}$ represent correlated Brownian motions, built on the filtered probability space $(\Omega, \mathcal{F}, \mathbb{F}, \mathbb{P})$ and such that
\[
\forall i,j \in \{1, \ldots, n\} \text{ and } \forall t > 0, \quad \operatorname{Cor}\left[B_t^{(i)}, B_t^{(j)}\right] = \rho_{ij}.
\]
There exists a matrix $\mathbf{A}$ of format $n \times n$ such that
\begin{itemize}
	\item[(i)] $\mathbf{B} = \mathbf{A}\mathbf{W}$
	\item[(ii)] $\operatorname{Cor}\left[B_t^{(i)}, B_t^{(j)}\right] = \rho_{ij}$
	\item[(iii)] $\mathbf{W}$ is made of independent Brownian motions.
\end{itemize}

\vspace{1em}

\textbf{Note:} $\mathbf{A}$ can be obtained by Cholesky decomposition: $\mathbf{A}= \frac{1}{\sqrt{t}}U^{\top}$ where $U$ is the upper triangular matrix coming from the Cholesky decomposition of the covariance matrix: $V_B=U^{\top}U$.

\vspace{1em}
Definition: For all $m \in \mathbb{N}$, the sequence of independent random variables $\{\xi_k^{(m)} : k \in \mathbb{N}\}$ is defined such that
\[
\xi_k^{(m)} = \begin{cases} -m^{-\frac{1}{2}} & \text{with probability } \frac{1}{2}, \\ m^{-\frac{1}{2}} & \text{with probability } \frac{1}{2}. \end{cases}
\]
Note that $\mathbb{E}\left[\xi_k^{(m)}\right] = 0$ and $\operatorname{Var}\left[\xi_k^{(m)}\right] = \frac{1}{m}$.
\vspace{1em}

Definition: \textbf{Random walk}
\vspace{0.5em}

For all $m \in \mathbb{N}$, let $X_0^{(m)} = 0$ and
\[
\forall n \in \left\{0, \frac{1}{m}, \frac{2}{m}, \ldots \right\}, \quad X_n^{(m)} = \sum_{k=1}^{mn} \xi_k^{(m)}.
\]
The process $X^{(m)} = \left\{X_n^{(m)} : n \in \left\{0, \frac{1}{m}, \frac{2}{m}, \ldots \right\}\right\}$ is a random walk.
\vspace{1em}

As $m$ increases, we take shorter and shorter steps (of length $m^{-\frac{1}{2}}$), as well as faster and faster (every $\frac{1}{m}$ time unit). Note that
\[
\mathbb{E}\left[X_n^{(m)}\right] = \sum_{k=1}^{mn} \mathbb{E}\left[\xi_k^{(m)}\right] = 0
\]
and
\[
\operatorname{Var}\left[X_n^{(m)}\right] = \sum_{k=1}^{mn} \operatorname{Var}\left[\xi_k^{(m)}\right] = \frac{mn}{m} = n.
\]
Moreover, the central limit theorem implies that $X_n^{(m)}$ converges in law to a $N(0, n)$ distribution as $m$ increases to infinity.
\vspace{1em}

\begin{itemize}
	\item Let's set
	\[
	Y_0^{(m)} = 0 \quad \text{and} \quad Y_t^{(m)} = X_{\lfloor n/m \rfloor}^{(m)} \quad \text{for all } t \in \left(\frac{n}{m}, \frac{n+1}{m}\right].
	\]
	The process $Y^{(m)}$ is defined for all $t \geq 0$. It is possible to show that this sequence $\left\{ Y^{(m)} : m \in \mathbb{N} \right\}$ of stochastic processes obtained from random walks converges in law to a Brownian motion process as $m$ tends to infinity. The challenge is to show that the limit has continuous paths.
\end{itemize}
\vspace{1em}

%-------------------
\section{Stochastic Integral}

Definition: \textbf{Basic stochastic process}
\vspace{0.5em}

We call $X$ a basic stochastic process if $X$ admits the following representation:
\[
X_t(\omega) = C(\omega) \mathbb{I}_{[a,b]}(t)
\]
where $a < b \in \mathbb{R}$ and $C$ is a random variable, $\mathcal{F}_a$--measurable and square-integrable, i.e. $\mathbb{E}^{\mathbb{P}}[C^2] < \infty$.
\vspace{1em}

Theorem: \textbf{Adapted and left-continuous processes are predictable}
\vspace{0.5em}

$\mathbb{F}$--adapted processes, the paths of which are left-continuous, are $\mathbb{F}$--predictable processes. In particular, $\mathbb{F}$--adapted processes with continuous paths are $\mathbb{F}$--predictable. (cf. Revuz and Yor)
\vspace{1em}

Definition: \textbf{Stochastic integral}
\vspace{0.5em}

The stochastic integral of $X$ with respect to the Brownian motion is defined as
\[
\left(\int_0^t X_s dW_s\right)(\omega) = C(\omega)(W_{t \wedge b}(\omega) - W_{t \wedge a}(\omega))
\]
\[
= \begin{cases}
0 & \text{if } 0 \leq t \leq a \\
C(\omega)(W_t(\omega) - W_a(\omega)) & \text{if } a < t < b \\
C(\omega)(W_b(\omega) - W_a(\omega)) & \text{if } b < t.
\end{cases}
\]
\vspace{1em}

Lemma: If $X$ is a basic stochastic process, then
\[
\left\{\int_0^t X_s dW_s : t \geq 0\right\}
\]
is a $(\Omega, \mathcal{F}, \mathbb{F}, \mathbb{P})$--martingale.
\vspace{1em}

Definition: \textbf{Stochastic integral with respect to a martingale}
\vspace{0.5em}

The stochastic integral of $X$ with respect to a square-integrable martingale $M$ is defined as
\[
\left(\int_0^t X_s dM_s\right)(\omega) = C(\omega)(M_{t \wedge b}(\omega) - M_{t \wedge a}(\omega))
\]
\[
= \begin{cases}
0 & \text{if } 0 \leq t \leq a \\
C(\omega)(M_t(\omega) - M_a(\omega)) & \text{if } a < t \leq b \\
C(\omega)(M_b(\omega) - M_a(\omega)) & \text{if } b < t.
\end{cases}
\]
\vspace{1em}

\textit{Note:} The proof of Lemma 1 is still valid if the Brownian motion is replaced with the martingale $M$. Therefore, $\left\{\int_0^t X_s dM_s\right\}_{t \geq 0}$ is a martingale.
\vspace{1em}

Lemma: \textbf{Expectation of product of integrals (Lemma 2)}
\vspace{0.5em}

If $X$ and $Y$ are basic processes, then for all $t \geq 0$,
\[
\mathbb{E}^{\mathbb{P}}\left[\int_0^t X_s Y_s \, ds\right] = \mathbb{E}^{\mathbb{P}}\left[\left(\int_0^t X_s dW_s\right)\left(\int_0^t Y_s dW_s\right)\right].
\]
\vspace{1em}

\textit{Also Note:} Since the Itô integral for basic processes is a martingale, then
\[
\mathbb{E}^{\mathbb{P}}\left[\int_0^t X_s dW_s\right] = \mathbb{E}^{\mathbb{P}}\left[\int_0^0 X_s dW_s\right] = 0.
\]
\vspace{1em}

Definition: \textbf{Simple stochastic process}
\vspace{0.5em}

We call $X$ a \textbf{simple stochastic process} if $X$ is a finite sum of basic processes:
\[
X_t(\omega) = \sum_{i=1}^{n} C_i(\omega)\mathbb{1}_{(a_i, b_i]}(t).
\]

where $a_1 < b_1 \leq a_2 < b_2 \leq \ldots \leq a_n < b_n$, and $C_i$ is $\mathcal{F}_{a_i}$--measurable.
\vspace{1em}

Definition: \textbf{Stochastic integral of X (simple processes)}
\vspace{0.5em}

The stochastic integral of $X$ with respect to the Brownian motion is defined as the sum of the stochastic integrals of the basic processes which constitute $X$:
\[
\int_0^t X_s dW_s = \sum_{i=1}^{n} \int_0^t C_i\mathbb{1}_{(a_i, b_i]}(s) dW_s.
\]
\vspace{1em}

Lemma: \textbf{Martingale property of stochastic integral (Lemma 3)}
\vspace{0.5em}

If $X$ is a simple stochastic process, then $\left\{\int_0^t X_s dW_s : t \geq 0\right\}$ is a $(\Omega, \mathcal{F}, \mathbb{F}, \mathbb{P})$--martingale.
\vspace{1em}

Lemma: \textbf{Itô isometry (Lemma 4)}
\vspace{0.5em}

If $X$ is a simple process, then for all $t \geq 0$,
\[
\mathbb{E}^{\mathbb{P}}\left[\int_0^t X_s^2 ds\right] = \mathbb{E}^{\mathbb{P}}\left[\left(\int_0^t X_s dW_s\right)^2\right].
\]
\vspace{1em}


\section{Stochastic differential equation and Ito’s lemma}

Theorem: \textbf{Itô's lemma}
\vspace{0.5em}

Let $W$ be a Brownian motion and let $f : \mathbb{R} \to \mathbb{R}$ be a function, the first two derivatives of which exist and are continuous. Then $\forall 0 \leq t \leq T$,
\[
f(W_t) - f(W_0) \stackrel{\mathbb{P}}{=a.s.} \int_0^t \frac{df}{dw}(W_s) dW_s + \frac{1}{2} \int_0^t \frac{d^2f}{dw^2}(W_s) ds. \tag{7}
\]
\vspace{1em}

In its differential form, Equation (7) is written
\[
df(W_t) = \frac{df}{dw}(W_t) dW_t + \frac{1}{2} \frac{d^2f}{dw^2}(W_t) dt.
\]
\vspace{1em}

Theorem: \textbf{Itô's lemma, first extension}
\vspace{0.5em}

Let $W$ be a Brownian motion and let $f : \mathbb{R} \times [0, \infty) \to \mathbb{R}$ be a function, the first and second partial derivatives of which exist and are continuous. Then $\forall 0 \leq t \leq T$,
\[
f(W_t, t) - f(W_0, 0) = \int_0^t \left(\frac{\partial f}{\partial t}(W_s, s) + \frac{1}{2} \frac{\partial^2 f}{\partial w^2}(W_s, s)\right) ds + \int_0^t \frac{\partial f}{\partial w}(W_s, s) dW_s.
\]

\textbf{In differential form:}
\[
df(W_t, t) = \left(\frac{\partial f}{\partial t}(W_t, t) + \frac{1}{2} \frac{\partial^2 f}{\partial w^2}(W_t, t)\right) dt + \frac{\partial f}{\partial w}(W_t, t) dW_t.
\]
\vspace{1em}

Definition: \textbf{Itô process}
\vspace{0.5em}

An \textbf{Itô process} is defined to be a process $X = \{X_t\}_{0 \leq t \leq T}$ such that:
\[
X_t \equiv X_0 + \int_0^t K_s ds + \int_0^t H_s dW_s \tag{8}
\]

where $K = \{K_t\}_{0 \leq t \leq T}$ and $H = \{H_t\}_{0 \leq t \leq T}$ are processes adapted to the Brownian filtration $\{\mathcal{F}_t\}_{0 \leq t \leq T}$, satisfying
\[
\mathbb{P}\left[\int_0^T |K_s| ds < \infty\right] = 1 \quad \text{and} \quad \mathbb{P}\left[\int_0^T |H_s|^2 ds < \infty\right] = 1.
\]
\vspace{1em}

Written in its differential form, the Itô process becomes
\[
dX_t = K_t dt + H_t dW_t.
\]
\vspace{1em}

Theorem: 

If 
\[
\mathbb{E}^{\mathbb{P}}\left[\int_0^T |H_s|^2 ds\right] < \infty, \tag{9}
\]

then the process $M = \{M_t\}_{0 \leq t \leq T}$ where $M_t = \int_0^t H_s dW_s$ is a $(\{\mathcal{F}_t\}_{0 \leq t \leq T}, \mathbb{P})$--martingale.
\vspace{1em}

Definition: \textbf{Quadratic variation}
\vspace{0.5em}

Let $0 = t_0^{(n)} \leq t_1^{(n)} \leq \ldots \leq t_n^{(n)} = T$ be a partition of the same interval $[0, T]$ such that $\lim_{n \to \infty} \max_{i \in \{1, \ldots, n\}} \left(t_i^{(n)} - t_{i-1}^{(n)}\right) = 0$. The \textbf{quadratic variation} of the stochastic process $X$ is
\[
\lim_{n \to \infty} \sum_{i=1}^{n} \left(X_{t_i^{(n)}} - X_{t_{i-1}^{(n)}}\right)^2
\]
\vspace{1em}

In the particular case where the stochastic process is an Itô process,
\[
X_t \equiv X_0 + \int_0^t K_s ds + \int_0^t H_s dW_s
\]

its \textbf{quadratic variation} is
\[
\langle X \rangle_t = \int_0^t H_s^2 ds.
\]
\vspace{1em}

Theorem: \textbf{Itô's lemma for Itô processes}
\vspace{0.5em}

Let $X$ be an Itô process and let $f : \mathbb{R} \to \mathbb{R}$ be a function, the first two derivatives of which exist and are continuous. Then $\forall 0 \leq t \leq T$,
\[
f(X_t) - f(X_0) \stackrel{\mathbb{P}}{=a.s.} \int_0^t \frac{df}{dX}(X_s) dX_s + \frac{1}{2} \int_0^t \frac{d^2f}{dX^2}(X_s) d\langle X \rangle_s \tag{10}
\]

\[
\stackrel{\mathbb{P}}{=a.s.} \int_0^t \frac{df}{dX}(X_s) H_s dW_s + \int_0^t \left(\frac{df}{dX}(X_s) K_s + \frac{1}{2} \frac{d^2f}{dX^2}(X_s) H_s^2 \right) ds,
\]

\textbf{In differential form:}
\[
df(X_t) = \frac{df}{dX}(X_t) dX_t + \frac{1}{2} \frac{d^2f}{dX^2}(X_t) d\langle X \rangle_t
= \left( \frac{df}{dX}(X_t) K_t + \frac{1}{2} \frac{d^2f}{dX^2}(X_t) H_t^2 \right) dt + \frac{df}{dX}(X_t) H_t dW_t
\]

\vspace{1em}

Theorem: \textbf{Fubini's theorem}
\vspace{0.5em}

If for all $0 \leq t \leq T$, $Y_t \geq 0$ then
\[
\mathbb{E}^{\mathbb{P}} \left[ \int_0^t Y_s ds \right] = \int_0^t \mathbb{E}^{\mathbb{P}} [Y_s] ds.
\]

\vspace{1em}

Definition: \textbf{Quadratic covariation}
\vspace{0.5em}

Let $X$ and $Y$ be two Itô processes such that
\[
dX_t = K_t dt + H_t dW_t \quad \text{and} \quad dY_t = \tilde{K}_t dt + \tilde{H}_t dW_t
\]

then the \textbf{quadratic covariation} process $\langle X, Y \rangle_t$ is defined for all $t \in [0, T]$ as
\[
\langle X, Y \rangle_t = \int_0^t H_s \tilde{H}_s ds.
\]

\vspace{1em}

Theorem: \textbf{The multiplication rule}
\vspace{0.5em}

Let $X$ and $Y$ be two Itô processes such that $dX_t = K_t dt + H_t dW_t$ and $dY_t = \tilde{K}_t dt + \tilde{H}_t dW_t$. The multiplication rule is
\[
dX_t Y_t = X_t dY_t + Y_t dX_t + d\langle X, Y \rangle_t
\]

\[
= X_t \left(\tilde{K}_t dt + \tilde{H}_t dW_t\right) + Y_t (K_t dt + H_t dW_t) + H_t \tilde{H}_t dt
\]

\[
= \left(X_t \tilde{K}_t + Y_t K_t + H_t \tilde{H}_t \right) dt + \left(X_t \tilde{H}_t + Y_t H_t \right) dW_t.
\]

\vspace{1em}

Definition: \textbf{Itô process (multidimensional)}
\vspace{0.5em}

Let $W^{(1)}, \ldots, W^{(n)}$ be independent Brownian motions. For $i, j \in \{1, \ldots, m\}$ and $k, k^* \in \{1, \ldots, n\}$, the adapted processes $K^{(i)}$ and $H^{(i,k)}$ satisfy
\[
\mathbb{P} \left[\sum_{i=1}^{m} \int_0^T |K_s^{(i)}| ds < \infty\right] = 1 \quad \text{and} \quad \mathbb{P} \left[\int_0^T H_s^{(i,k)} H_s^{(j,k^*)} ds < \infty\right] = 1.
\]

For all $i \in \{1, 2, \ldots, m\}$, we define the process $X^{(i)}$ as
\[
X_t^{(i)} = X_0^{(i)} + \int_0^t K_s^{(i)} ds + \sum_{k=1}^{n} \int_0^t H_s^{(i,k)} dW_s^{(k)} \tag{12}
\]

where $X_0^{(1)}, \ldots, X_0^{(m)}$ are $\mathcal{F}_0$--measurable random variables.

\textbf{In differential form:}
\[
dX_t^{(i)} = K_t^{(i)} dt + \sum_{k=1}^{n} H_t^{(i,k)} dW_t^{(k)}
\]

\vspace{1em}

Definition: \textbf{Quadratic covariation (multidimensional)}
\vspace{0.5em}

For two Itô processes $X^{(i)}$ and $X^{(j)}$, the \textbf{quadratic covariation} is defined for all $t \in [0, T]$ as
\[
\langle X^{(i)}, X^{(j)} \rangle_t = \sum_{k=1}^{n} \int_0^t H_s^{(i,k)} H_s^{(j,k)} ds.
\]

\textbf{Note that it is symmetric,}
\[
\langle X^{(i)}, X^{(j)} \rangle = \langle X^{(j)}, X^{(i)} \rangle,
\]

and the quadratic variation is a particular case
\[
\langle X^{(i)}, X^{(i)} \rangle = \langle X^{(i)} \rangle.
\]

\vspace{1em}

Theorem: \textbf{Itô's lemma, multivariate version}
\vspace{0.5em}

Let $f : \mathbb{R}^m \times [0, \infty) \to \mathbb{R}$ be a continuous function such that the partial derivatives $\frac{\partial f}{\partial x_i}$, $\frac{\partial f}{\partial s}$, $\frac{\partial^2 f}{\partial x_i \partial x_j}$, and $\frac{\partial^2 f}{\partial x_i \partial s}$ for $i,j \in \{1, \ldots, m\}$ exist and are continuous. Then $\forall t \geq 0$,
\[
f(X_t^{(1)}, \ldots, X_t^{(m)}, t) - f(X_0^{(1)}, \ldots, X_0^{(m)}, 0)
\]

\[
\stackrel{\mathbb{P}}{=a.s.} \sum_{i=1}^{m} \int_0^t \frac{\partial f}{\partial x_i}(X_s^{(1)}, \ldots, X_s^{(m)}, s) dX_s^{(i)} + \int_0^t \frac{\partial f}{\partial t}(X_s^{(1)}, \ldots, X_s^{(m)}, s) ds
\]

\[
+ \frac{1}{2} \sum_{i=1}^{m} \sum_{j=1}^{m} \int_0^t \frac{\partial^2 f}{\partial x_i \partial x_j}(X_s^{(1)}, \ldots, X_s^{(m)}, s) d\langle X^{(i)}, X^{(j)} \rangle_s.
\]

\textbf{In differential form:}
\[
df(X_t^{(1)}, \ldots, X_t^{(m)}, t) = \sum_{i=1}^{m} \frac{\partial f}{\partial x_i}(X_t^{(1)}, \ldots, X_t^{(m)}, t) dX_t^{(i)} + \frac{\partial f}{\partial t}(X_t^{(1)}, \ldots, X_t^{(m)}, t) dt
\]

\[
+ \frac{1}{2} \sum_{i=1}^{m} \sum_{j=1}^{m} \frac{\partial^2 f}{\partial x_i \partial x_j}(X_t^{(1)}, \ldots, X_t^{(m)}, t) d\langle X^{(i)}, X^{(j)} \rangle_t.
\]

\vspace{1em}

Definition: \textbf{Strong solution}
\vspace{0.5em}

A \textbf{strong solution} consists of finding a stochastic process $X$ existing on the same filtered probability space $(\Omega, \mathcal{F}, \{\mathcal{F}_t : t \geq 0\}, \mathbb{P})$ as the Brownian motion and satisfying the equation
\[
X_t = X_0 + \int_0^t \mu(X_s) ds + \int_0^t \sigma(X_s) dW_s.
\]

\textbf{Uniqueness:} The strong solution is said to be \textbf{unique} if, when $X$ and $\tilde{X}$ represent two strong solutions to the same stochastic differential equation, we have
\[
\mathbb{P}[X_t = \tilde{X}_t, \, t \geq 0] = 1.
\]

\vspace{1em}

Definition: \textbf{Weak solution}
\vspace{0.5em}

We have a \textbf{weak solution} if we can construct a probability space $(\tilde{\Omega}, \tilde{\mathcal{F}}, \tilde{\mathbb{P}})$, a $(\Omega, \mathcal{F}, \mathbb{P})$ – standard Brownian motion $W$, and a $(\tilde{\Omega}, \tilde{\mathcal{F}}, \tilde{\mathbb{P}})$ – stochastic process $\tilde{X}$ such that
\[
\tilde{X}_t = \tilde{X}_0 + \int_0^t \mu(\tilde{X}_s) ds + \int_0^t \sigma(\tilde{X}_s) d\tilde{W}_s.
\]

\textbf{Relationship with strong solutions:} It is possible to have several strong solutions to the same equation. If the strong solution is unique, then there will also be a unique weak solution. By contrast, if the weak solution is unique, there may be several strong solutions.

\vspace{1em}

Definition: \textbf{Gaussian models}
\vspace{0.5em}

Some stochastic processes that give rise to \textbf{Gaussian models} are of the form
\[
dX_t = (\mu_1(t) X_t + \mu_2(t)) dt + \sigma(t) dW_t \tag{13}
\]

where $\mu_1(\cdot)$, $\mu_2(\cdot)$ and $\sigma(\cdot)$ are deterministic functions of time and $\{W_t\}_{0 \leq t \leq T}$ is a $\mathbb{P}$--standard Brownian motion.

\vspace{1em}

Theorem: \textbf{Strong solution to Gaussian models}
\vspace{0.5em}

If $\mathbb{P} \left[\int_0^T \frac{\sigma^2(s)}{\Phi^2(s)} ds < \infty\right] = 1$, then the strong solution to the stochastic differential equation (13) is
\[
X_t = \Phi(t) \left[X_0 + \int_0^t \frac{\mu_2(s)}{\Phi(s)} ds + \int_0^t \frac{\sigma(s)}{\Phi(s)} dW_s\right]
\]

where the function $\Phi(\cdot)$ is the solution to the ordinary differential equation
\[
d\Phi(t) = \mu_1(t) \Phi(t) dt, \quad \Phi(0) = 1.
\]

\textit{Source: Bisière (1997), page 117.}

\vspace{1em}

Theorem: \textbf{Lipschitz condition for existence of strong solutions}
\vspace{0.5em}

If for all $i, j, x$, and $y$, we have $|\sigma_{i,j}(x) - \sigma_{i,j}(y)| \leq K|x - y|$ and $|b_i(x) - b_i(y)| \leq K|x - y|$, then the integral equation
\[
X_t = x + \int_0^t b(X_s) ds + \int_0^t \sigma(X_s) dW_s
\]

has a strong solution.


\vspace{1em}


\section{Girsanov's theorem}

Definition: \textbf{Equivalent probability measures}
\vspace{0.5em}

Two probability measures $\mathbb{P}$ and $\mathbb{Q}$ constructed on the same measurable space $(\Omega, \mathcal{F})$ are said to be \textbf{equivalent} if they have the same set of impossible events, i.e.

\[
\mathbb{P}(A) = 0 \Leftrightarrow \mathbb{Q}(A) = 0, \quad A \in \mathcal{F}.
\]

\vspace{1em}

Theorem: \textbf{Radon-Nikodym theorem}
\vspace{0.5em}

Given two equivalent probability measures $\mathbb{P}$ and $\mathbb{Q}$ constructed on the measurable space $(\Omega, \mathcal{F})$, there exists a positive-valued random variable $Y$ such that
\[
\mathbb{Q}(A) = \mathbb{E}^{\mathbb{P}}[Y\mathbb{I}_A].
\]

Such a random variable $Y$ is often denoted by $\frac{d\mathbb{Q}}{d\mathbb{P}}$.

\vspace{1em}


Theorem: \textbf{Cameron-Martin-Girsanov}
\vspace{0.5em}

Let $\gamma = \{\gamma_t : t \in [0, T]\}$ be a $\{\mathcal{F}_t\}$--predictable process such that $\mathbb{E}^{\mathbb{P}}\left[\exp\left(\frac{1}{2}\int_0^T \gamma_t^2 dt\right)\right] < \infty$. There exists a measure $\mathbb{Q}$ on $(\Omega, \mathcal{F})$ such that:

\begin{itemize}
	\item (G1) $\mathbb{Q}$ is equivalent to $\mathbb{P}$,
	\item (G2) $\displaystyle \frac{d\mathbb{Q}}{d\mathbb{P}} = \exp\left[-\int_0^T \gamma_t dW_t^{\mathbb{P}} - \frac{1}{2}\int_0^T \gamma_t^2 dt\right]$,
	\item (G3) The process $\{W_t^{\mathbb{Q}}\}_{0 \leq t \leq T}$ defined as $W_t^{\mathbb{Q}} = W_t^{\mathbb{P}} + \int_0^t \gamma_s ds$ is a $(\{\mathcal{F}_t\}, \mathbb{Q})$--Brownian motion.
\end{itemize}

\vspace{0.5em}

\textit{Condition:} The condition $\mathbb{E}^{\mathbb{P}}\left[\exp\left(\frac{1}{2}\int_0^T \gamma_t^2 dt\right)\right] < \infty$ is sufficient but not necessary. It is known as the \textbf{Novikov condition}.

\vspace{1em}


Theorem: \textbf{Theorem (Cameron-Martin-Girsanov theorem) - Multivariate case}
\vspace{0.5em}

Let $W = (W^{\mathbb{P}, 1}, \ldots, W^{\mathbb{P}, n})$ be $n$ independent $(\{\mathcal{F}_t\}, \mathbb{P})$-standard Brownian motions. For all $i \in \{1, 2, \ldots, n\}$, assume that $\gamma^{(i)}$ is a $\{\mathcal{F}_t\}$--predictable process such that

\[
\mathbb{E}^{\mathbb{P}}\left[\exp\left(\frac{1}{2}\int_0^T \left(\gamma_t^{(i)}\right)^2 dt\right)\right] < \infty.
\]

There exists a measure $\mathbb{Q}$ on $(\Omega, \mathcal{F})$ such that:

\begin{itemize}
	\item (G1) $\mathbb{Q}$ is equivalent to $\mathbb{P}$,
	\item (G2) $\displaystyle \frac{d\mathbb{Q}}{d\mathbb{P}} = \exp\left[-\sum_{i=1}^n \int_0^T \gamma_t^{(i)} dW_t^{\mathbb{P}, i} - \frac{1}{2}\int_0^T \sum_{i=1}^n \left(\gamma_t^{(i)}\right)^2 dt\right]$,
	\item (G3) For all $i \in \{1, \ldots, n\}$, the process $W_t^{\mathbb{Q}, i}$ defined as 
	\[
	W_t^{\mathbb{Q}, i} = W_t^{\mathbb{P}, i} + \int_0^t \gamma_s^{(i)} ds
	\]
	is a $(\{\mathcal{F}_t\}, \mathbb{Q})$--Brownian motion.
\end{itemize}

\vspace{1em}


\section{Hedging attainable contingent claims}

Definition: \textbf{Trading strategy}
\vspace{0.5em}

A trading strategy is a predictable portfolio process
\[
\{(\phi_t, \psi_t) : 0 \leq t \leq T\}
\]
where
\begin{itemize}
	\item $\phi_t$ = the number of riskless asset shares held at time $t$
	\item $\psi_t$ = the number of risky asset shares held at time $t$
\end{itemize}
\vspace{1em}

The value of such a trading strategy at time $t$ is
\[
V_t = \phi_t B_t + \psi_t S_t.
\]
\[
dV_t = \phi_t dB_t + B_t d\phi_t + d\langle\phi, B\rangle_t + \psi_t dS_t + S_t d\psi_t + d\langle\psi, S\rangle_t.
\]
\vspace{1em}

Definition: \textbf{Self-financing strategy}
\vspace{0.5em}

A trading strategy is said to be self-financing if
\[
dV_t = \phi_t dB_t + \psi_t dS_t,
\]
i.e. the variations in the strategy value are only caused by the asset price variations and do not depend on the portfolio weight fluctuations. A strategy is therefore self-financing if
\[
B_t d\phi_t + S_t d\psi_t + d \langle\phi, B\rangle_t + d \langle\psi, S\rangle_t = 0.
\]
\vspace{1em}

Definition: \textbf{Attainable contingent claim}
\vspace{0.5em}

Let $X$ be a contingent claim. $X$ is therefore a non-negative-valued, $\mathcal{F}_T$--measurable, random variable. The contingent claim $X$ is attainable if there exists a self-financing trading strategy $\{(\phi_t, \psi_t) : 0 \leq t \leq T\}$ such that
\[
V_T = \phi_T B_T + \psi_T S_T = X.
\]
\vspace{1em}

\end{document}
